% This file contains the templates for the first few pages of the thesis including 
% 1. Title page
% 2. Dedication
% 3. Dissertation approval
% 4. Declaration of authorship
% 5. Abstract


%   1. TITLE
\newcommand{\titlePage}{

% No page number
\thispagestyle{empty}
\begin{center}

% thesis title
\vspace*{15px}
{\Huge\bfseries \thesisTitle}\\[1.0cm] 

% submitted in partial fulfillment etc.
\textit{Submitted in partial fulfillment of the requirements\\[0.2cm] of the degree of}\\[0.2cm] \degree\\
in \\ 
Communication and Signal Processing \\[1.0cm]
% author
\textit{by}\\[0.1cm]
\authorName \\[0.1cm] (\textit{Roll no.} \rollNo) \\[2.0cm]

% supervisor
\textit{Supervisor:}\\[0.2cm]
% or
% \textit{under the supervision of}\\[0.2cm]
\supervisorOne \\[1.0cm]

% iit-b logo
\includegraphics[width=0.25\textwidth]{Figures/iitb_logo.jpg}

\vspace*{10px}

% department and college
\dept\\[0.2cm]
\college\\[0.2cm]

% year
\currentmonth \ \currentyear\\ 

\end{center}

\clearpage
}



%   3. DISSERTATION APPROVAL

\newcommand{\approval}{
% if you use a dedication, then use a page number with the plain page style
\thispagestyle{plain}
% else, use no page number with the empty page style
%\thispagestyle{empty}


% page title
\begin{center}{\huge\bf M.Tech Dissertation Approval\par}\end{center}

\vspace*{15px}

\noindent This dissertation entitled 
% title
\textbf{``\thesisTitle"}, submitted by 
% author
\authorName  
 \ (Roll No.  \rollNo),
is approved for the award of degree of 
% degree
\degree \
in 
% branch
Communication and signal processing\\[1.0cm]

% examiners and supervisors
\begin{flushright}
\textbf{{Examiners}}\\[0.8cm]
\examinerOne\quad\rule{0.3\textwidth}{.3pt}\\[0.8cm]
\examinerTwo\quad\rule{0.3\textwidth}{.3pt}\\[1.6cm]

\textbf{{Supervisor}}\\[0.8cm]
\quad\rule{0.3\textwidth}{.3pt}\\[1.6cm]

\textbf{{Chairman}}\\[0.8cm]
\chairman\quad\rule{0.3\textwidth}{.3pt}\\[2.4cm]

\end{flushright}
\textbf{Date:} 05 \currentmonth { } \currentyear\\[0.3cm]
\textbf{Place:} \  Mumbai

% start a new page
\clearpage
}

%   4. DECLARATION OF AUTHORSHIP
\newcommand{\authorship}{
\thispagestyle{plain}

\begin{center}{\huge\bf Declaration of Authorship\par}\end{center}

\vspace*{15px}

% institute declaration text
\noindent I declare that this written submission represents my ideas in my own words and where others' ideas or words have been included, I have adequately cited and referenced the original sources.  I also declare that I have adhered to all principles of academic honesty and integrity and   have   not   misrepresented   or   fabricated   or   falsified   any   idea/data/fact/source   in   my submission.  I understand that any violation of the above will be cause for disciplinary action by the Institute and can also evoke  penal action from the sources which have thus not been properly cited or from whom proper permission has not been taken when needed.

\vspace*{10px}

% signature
\begin{flushright}
{Signature: ......................................\\[0.4cm]}

% author name
{\textbf{\authorName}\\[0.0cm]\rollNo\\[2.0cm]}

\end{flushright}
% date
\begin{flushleft}
{Date: 05 \currentmonth { } \currentyear\\}
\end{flushleft}


% start a new page
\clearpage 
}

%   5. ABSTRACT
\newcommand{\abstractpage}{
\thispagestyle{plain}

% header-type text for the abstract - optional
\small {\noindent\authorName/ \supervisorOne{ }(Supervisor): \textbf{``\thesisTitle"}, \textit{M.Tech Dissertation}, \dept, \college, \currentmonth { } \currentyear.}\\[0.0cm]
\HRule\\[0.2cm]


\vspace*{10px}

\begin{center}{\huge{\textit {Abstract}}\par}\end{center}

\vspace*{10px}

% Abstract text - Type abstract here
\noindent Dereverberation of recorded speech signals is one of the most pertinent problems in speech processing. In the present work, the objective is to understand and implement dereverberation techniques that aim at enhancing the magnitude spectrogram of reverberant speech signals to remove the reverberant effects introduced. An approach to estimate clean speech spectrogram from the reverberant speech spectrogram is proposed. This is achieved through non-negative matrix factor deconvolution(NMFD). Further, this approach is extended using the NMF representation for speech magnitude spectrograms. To exploit temporal dependencies, a convolutive NMF based representation and a frame stacked model are incorporated into the  NMFD framework for speech.
A novel approach for dereverberation by applying NMFD to the activation matrix of the reverberated magnitude spectrogram is also proposed .
Finally, a comparative analysis of the performance of the listed techniques, using sentence recordings from the TIMIT database and recorded room impulse responses from the Reverb 2014 challenge are presented based on two key objective measures - PESQ and Cepstral Distortion.\\
Although, we were qualitatively able to verify the claims made in literature regarding these techniques, exact results could not be matched. The novel approach, as it is suggested, provides improvement in quantitative metrics, but is not consistent



% Start a new page
\clearpage 
}
